\documentclass[../../main.tex]{subfiles}
\begin{document}

\begin{center}
  {\fontsize{48pt}{58pt}\selectfont \textbf{1987}\normalsize\textbf{年度}}
\end{center}

\begin{table}[htb]
  \centering
  \caption{出題テーマと難易度}
  \begin{tabular}{|c|p{6.2cm}|c|c|c|}
    \hline
    \textbf{問題番号} & \textbf{テーマ} & \textbf{思考量} & \textbf{計算量} & \textbf{総合} \\
    \hline
    第1問 & 一次変換による直線の像と角度 &  &  &  \\
    \hline
    第2問 & 平行移動した放物線と双曲線の接点 & 普 & 難 & 普 \\
    \hline
    第3問 & 球の接平面と三角形の面積 & 難 & 難 & 難 \\
    \hline
    第4問 & 回転してできる立体の体積 & 普 & 普 & 普 \\
    \hline
    第5問 & 数列の並べ替えと平方差の和 & 普 & 普 & 普 \\
    \hline
    第6問 & 正六角形上の直角三角形と確率 & 難 & 難 & 難 \\
    \hline
  \end{tabular}
\end{table}

\begin{table}[htb]
  \centering
  \caption{解答}
  \renewcommand{\arraystretch}{1.4}
  \begin{tabular}{|c|c|p{9cm}|}
    \hline
    \textbf{問題番号} & \textbf{小問} & \textbf{答え} \\
    \hline
    第1問 & -- & \\
    \hline
    第2問 & -- & 点$(a,b)$の領域は解答の図を参照．他に共有点がないのは$\left(\dfrac{3}{2},\dfrac{3}{4}\right)$ \\
    \hline
    第3問 & -- & 面積の最小値は$3+2\sqrt{2}$ \\
    \hline
    第4問 & -- & $V=\dfrac{\pi}{3}(y^2-1)^2+\dfrac{2\pi}{3}$，最小値は$\dfrac{2\pi}{3}$ \\
    \hline
    第5問 & -- & 与えられた不等式が成立（証明問題） \\
    \hline
    第6問 & (1) & $p_3=\dfrac{1}{3}$，$p_4=\dfrac{11}{18}$ \\
    \cline{2-3}
    & (2) & 極限値は$-\log 2$ \\
    \hline
  \end{tabular}
\end{table}

\paragraph{試験形態}
6問．試験時間は要確認．

\end{document}
